%% toc: P2 17 17.4,17.5,17.6
\section{Proofs for first-order Hilbert deduction and Kleene's deduction theorem}
\label{sec:annex-p02-w02-p02-025-fol-calculi}

Notation is that of \cref{sec:p02-w02-p02-025-fol-calculi}.  A line of a
dependency-annotated derivation is a pair \((\chi,D)\) together with its
justification: premise, axiom-schema instance, or the rule and the earlier
lines it is applied to.  For formulas \(\alpha,\beta\) of the fragment,
\(\FV(\alpha\to\beta)=\FV(\alpha)\cup\FV(\beta)\) by the definition of free
variables.

\begin{annexitem}{annex:p02-w02-p02-025-identity}{thm:p02-w02-p02-025-identity}{Identity with empty support}
\begin{proof}
Line~1 is the A1 instance \(\varphi\to(\psi\to\varphi)\) with
\(\psi\defeq\varphi\to\varphi\).  Line~2 is the A2 instance
\((\varphi\to(\psi\to\chi))\to((\varphi\to\psi)\to(\varphi\to\chi))\) with
\(\psi\defeq\varphi\to\varphi\) and \(\chi\defeq\varphi\).  Line~2 has the
form \(\alpha\to\beta\) with \(\alpha\) equal to line~1, so modus ponens gives
\(\beta\), which is line~3.  Line~4 is the A1 instance with
\(\psi\defeq\varphi\).  Line~3 has the form \(\alpha\to\beta\) with \(\alpha\)
equal to line~4, and modus ponens gives line~5, \(\varphi\to\varphi\).
Axiom lines have empty support and modus ponens takes the union of the
supports of its premises, so every support is \(\varnothing\).  The schemas
A1 and A2 of \texttt{sheet:\allowbreak hilbert-fol} admit arbitrary fragment
formulas for their metavariables, so the derivation exists for every
fragment formula \(\varphi\).
\end{proof}
\end{annexitem}

\begin{annexitem}{annex:p02-w02-p02-025-permutation}{thm:p02-w02-p02-025-permutation}{Permutation of antecedents with empty support}
\begin{proof}
Fix \(\varphi,\theta,\psi\) and abbreviate
\[
 \alpha\defeq\varphi\to(\theta\to\psi),\quad
 \beta\defeq\varphi\to\theta,\quad
 \gamma\defeq\varphi\to\psi,
\]
\[
 \kappa\defeq\theta\to\beta,\quad
 \delta\defeq\theta\to\gamma,\quad
 \sigma\defeq\beta\to\gamma,\qquad
 \pi\defeq\theta\to\sigma,\qquad
 \rho\defeq\kappa\to\delta .
\]
The target is \(\alpha\to\delta\), which is
\((\varphi\to(\theta\to\psi))\to(\theta\to(\varphi\to\psi))\).  The
following nineteen lines form a derivation; every line has empty support,
since only axiom lines and modus ponens occur.
\[
\begin{array}{lll}
 1. & \alpha\to\sigma & \text{A2}\\
 2. & \pi\to\rho & \text{A2}\\
 3. & (\pi\to\rho)\to(\sigma\to(\pi\to\rho)) & \text{A1}\\
 4. & \sigma\to(\pi\to\rho) & \text{MP }3,2\\
 5. & (\sigma\to(\pi\to\rho))\to((\sigma\to\pi)\to(\sigma\to\rho)) & \text{A2}\\
 6. & (\sigma\to\pi)\to(\sigma\to\rho) & \text{MP }5,4\\
 7. & \sigma\to\pi & \text{A1}\\
 8. & \sigma\to\rho & \text{MP }6,7\\
 9. & (\sigma\to\rho)\to(\alpha\to(\sigma\to\rho)) & \text{A1}\\
 10. & \alpha\to(\sigma\to\rho) & \text{MP }9,8\\
 11. & (\alpha\to(\sigma\to\rho))\to((\alpha\to\sigma)\to(\alpha\to\rho)) & \text{A2}\\
 12. & (\alpha\to\sigma)\to(\alpha\to\rho) & \text{MP }11,10\\
 13. & \alpha\to\rho & \text{MP }12,1\\
 14. & (\alpha\to\rho)\to((\alpha\to\kappa)\to(\alpha\to\delta)) & \text{A2}\\
 15. & (\alpha\to\kappa)\to(\alpha\to\delta) & \text{MP }14,13\\
 16. & \kappa & \text{A1}\\
 17. & \kappa\to(\alpha\to\kappa) & \text{A1}\\
 18. & \alpha\to\kappa & \text{MP }17,16\\
 19. & \alpha\to\delta & \text{MP }15,18
\end{array}
\]
The schema instances are checked against A1
\(\varphi'\to(\psi'\to\varphi')\) and A2
\((\varphi'\to(\psi'\to\chi'))\to((\varphi'\to\psi')\to(\varphi'\to\chi'))\),
with primed metavariables to avoid a clash with the fixed formulas.
\begin{itemize}
  \item Line~1: A2 with \(\varphi'\defeq\varphi\), \(\psi'\defeq\theta\),
        \(\chi'\defeq\psi\); its antecedent is \(\alpha\) and its consequent
        is \((\varphi\to\theta)\to(\varphi\to\psi)=\beta\to\gamma=\sigma\).
  \item Line~2: A2 with \(\varphi'\defeq\theta\), \(\psi'\defeq\beta\),
        \(\chi'\defeq\gamma\); its antecedent is
        \(\theta\to(\beta\to\gamma)=\pi\) and its consequent is
        \((\theta\to\beta)\to(\theta\to\gamma)=\kappa\to\delta=\rho\).
  \item Line~3: A1 with \(\varphi'\defeq\pi\to\rho\),
        \(\psi'\defeq\sigma\).
  \item Line~5: A2 with \(\varphi'\defeq\sigma\), \(\psi'\defeq\pi\),
        \(\chi'\defeq\rho\).
  \item Line~7: \(\sigma\to\pi\) is \(\sigma\to(\theta\to\sigma)\), which is
        A1 with \(\varphi'\defeq\sigma\), \(\psi'\defeq\theta\).
  \item Line~9: A1 with \(\varphi'\defeq\sigma\to\rho\),
        \(\psi'\defeq\alpha\).
  \item Line~11: A2 with \(\varphi'\defeq\alpha\), \(\psi'\defeq\sigma\),
        \(\chi'\defeq\rho\).
  \item Line~14: \(\alpha\to\rho\) is \(\alpha\to(\kappa\to\delta)\), so the
        line is A2 with \(\varphi'\defeq\alpha\), \(\psi'\defeq\kappa\),
        \(\chi'\defeq\delta\).
  \item Line~16: \(\kappa=\theta\to(\varphi\to\theta)\) is A1 with
        \(\varphi'\defeq\theta\), \(\psi'\defeq\varphi\).
  \item Line~17: A1 with \(\varphi'\defeq\kappa\), \(\psi'\defeq\alpha\).
\end{itemize}
Each modus ponens step ``MP \(m,k\)'' has line \(m\) of the form
\(\lambda\to\lambda'\) with \(\lambda\) literally equal to line \(k\), and its
conclusion is \(\lambda'\): line~3 has antecedent line~2; line~5 has
antecedent line~4; line~6 has antecedent line~7; line~9 has antecedent
line~8; line~11 has antecedent line~10; line~12 has antecedent line~1;
line~14 has antecedent line~13; line~17 has antecedent line~16; and line~15
has antecedent line~18.  Expanding the abbreviations, line~19 is
\((\varphi\to(\theta\to\psi))\to(\theta\to(\varphi\to\psi))\).

For the consequence, let a derivation from \(\Gamma\) end in
\((\varphi\to(\theta\to\psi),E)\).  Append the nineteen lines, each with
empty support, and then apply modus ponens to line~19 and the line
\((\varphi\to(\theta\to\psi),E)\).  The conclusion is
\(\theta\to(\varphi\to\psi)\) with support \(\varnothing\cup E=E\).  The
appended lines are axiom instances and modus ponens steps, which are legal
in a derivation from any premise set.  The inverse passage is the same
construction for the instance obtained by putting \(\theta\) for
\(\varphi\) and \(\varphi\) for \(\theta\), namely
\((\theta\to(\varphi\to\psi))\to(\varphi\to(\theta\to\psi))\).
\end{proof}
\end{annexitem}

\begin{annexitem}{annex:p02-w02-p02-025-support}{thm:p02-w02-p02-025-support}{Concatenation and support restriction}
\begin{proof}
\emph{Clause 1.}  By \cref{def:p02-hilbert-fol-dependencies}, whether a line
is correctly justified depends only on its own formula and support and on
the formulas and supports of the earlier lines it cites: a premise line
needs its formula in \(\Gamma\) and its support equal to the singleton; an
axiom line needs an instance of a schema and the empty support; modus ponens
needs the two cited lines to have the forms \(\eta\) and \(\eta\to\chi\) and
the support to be their union; generalization and existential elimination
need the cited line to have the displayed form, the same support, and the
variable conditions, which mention only that line's formula and support.
In \(\mathcal D_1\mathcal D_2\) each line of \(\mathcal D_1\) keeps its
position and citations, and each line of \(\mathcal D_2\) cites the same
lines as before, shifted by the length of \(\mathcal D_1\).  So every
justification remains correct, and every line keeps its support.

\emph{Clause 2.}  Let the line \((\chi,D)\) be line \(i\) of a derivation
\[
 \langle(\chi_1,D_1),\ldots,(\chi_n,D_n)\rangle
\]
from \(\Gamma\).  Let
\(S\) be the least set of indices containing \(i\) and containing every index
\(j\) such that some \(k\in S\) uses line \(j\).  Every member of \(S\) is at
most \(i\), because a line uses only earlier lines.

First, \(D_j\subseteq D_i\) for every \(j\in S\).  If line \(k\) uses line
\(j\), then \(D_j\subseteq D_k\): modus ponens gives \(D_k\) as the union of
the supports of the two lines it uses, and generalization and existential
elimination give \(D_k=D_j\).  Every \(j\in S\) is reached from \(i\) by a
finite chain of uses, so induction on the length of the chain gives
\(D_j\subseteq D_i=D\subseteq\Delta\).

Now list the lines with indices in \(S\) in increasing order, keeping each
line's formula, support and justification, and renumbering citations.
Each listed line is justified as before.  A premise line \((\chi_j,\{\chi_j\})\)
has \(\chi_j\in D_j\subseteq\Delta\), so it is a premise line of a
derivation from \(\Delta\).  An axiom line needs nothing more.  A rule line
cites only lines it uses, which belong to \(S\) by the definition of \(S\)
and precede it in the list, and its formula, support and variable conditions
are unchanged.  So the list is a derivation from \(\Delta\), and its last
line is line \(i\), \((\chi,D)\), because \(i\) is the greatest member of
\(S\).
\end{proof}
\end{annexitem}

\begin{annexitem}{annex:p02-w02-p02-025-discharge}{thm:p02-w02-p02-025-discharge}{Hilbert-FOL deduction theorem with explicit supports}
\begin{proof}
\emph{Clause 2.}  Let \(\mathcal D\) be a derivation from \(\Gamma\) ending in
\((\varphi\to\psi,E)\).  It is also a derivation from
\(\Gamma\cup\{\varphi\}\): its premise lines have formulas in
\(\Gamma\subseteq\Gamma\cup\{\varphi\}\), and no other justification refers to
the premise set.  Append the premise line \((\varphi,\{\varphi\})\) and apply
modus ponens to \((\varphi\to\psi,E)\) and \((\varphi,\{\varphi\})\).  The
result is \((\psi,E\cup\{\varphi\})\).

\emph{Clause 1.}  Let \(\langle(\chi_1,D_1),\ldots,(\chi_n,D_n)\rangle\) be a
derivation from \(\Gamma\cup\{\varphi\}\) with \(\chi_n=\psi\) and
\(D_n=D\).  By strong induction on \(i\leq n\) we construct a derivation
\(\mathcal D_i\) from \(\Gamma\) whose last line is
\((\varphi\to\chi_i,E_i)\) with \(E_i\subseteq D_i\setminus\{\varphi\}\).
For \(i=n\) this is the claim.  Every \(\mathcal D_i\) below is a
concatenation of earlier \(\mathcal D_j\), derivations from \(\Gamma\), and
new lines whose justifications are written out; by
\cref{thm:p02-w02-p02-025-support}(1) it is a derivation from \(\Gamma\).
There are six cases, according to the justification of line \(i\).

\emph{Case (a): line \(i\) is a premise line and \(\chi_i=\varphi\).}  Let
\(\mathcal D_i\) be the five-line derivation of
\cref{thm:p02-w02-p02-025-identity}.  It ends in
\((\varphi\to\varphi,\varnothing)\), and \(E_i=\varnothing\).

\emph{Case (b): line \(i\) is a premise line and \(\chi_i\neq\varphi\).}
Then \(\chi_i\in\Gamma\) and \(D_i=\{\chi_i\}\).  Let \(\mathcal D_i\) be
\[
 (\chi_i,\{\chi_i\})\ \text{premise},\qquad
 (\chi_i\to(\varphi\to\chi_i),\varnothing)\ \text{A1},\qquad
 (\varphi\to\chi_i,\{\chi_i\})\ \text{MP}.
\]
Here \(E_i=\{\chi_i\}=D_i\setminus\{\varphi\}\), since
\(\chi_i\neq\varphi\).

\emph{Case (c): line \(i\) is an axiom-schema instance.}  Then
\(D_i=\varnothing\).  Let \(\mathcal D_i\) be
\[
 (\chi_i,\varnothing)\ \text{axiom},\qquad
 (\chi_i\to(\varphi\to\chi_i),\varnothing)\ \text{A1},\qquad
 (\varphi\to\chi_i,\varnothing)\ \text{MP},
\]
so \(E_i=\varnothing\).

\emph{Case (d): line \(i\) comes by modus ponens from lines \(j,k<i\)} with
\(\chi_j=\eta\) and \(\chi_k=\eta\to\chi_i\), so \(D_i=D_j\cup D_k\).  By the
induction hypothesis \(\mathcal D_j\) ends in \((\varphi\to\eta,E_j)\) and
\(\mathcal D_k\) ends in \((\varphi\to(\eta\to\chi_i),E_k)\).  Let
\(\mathcal D_i\) be \(\mathcal D_j\mathcal D_k\) followed by
\[
\begin{array}{ll}
 ((\varphi\to(\eta\to\chi_i))\to((\varphi\to\eta)\to(\varphi\to\chi_i)),
   \varnothing) & \text{A2},\\
 ((\varphi\to\eta)\to(\varphi\to\chi_i),\ E_k) & \text{MP, last line of }\mathcal D_k,\\
 (\varphi\to\chi_i,\ E_j\cup E_k) & \text{MP, last line of }\mathcal D_j.
\end{array}
\]
Then
\(E_i=E_j\cup E_k\subseteq(D_j\setminus\{\varphi\})\cup(D_k\setminus\{\varphi\})
=D_i\setminus\{\varphi\}\).

\emph{Case (e): line \(i\) comes by generalization from line \(j<i\).}
Then \(\chi_i=\forall x\,\chi_j\), \(D_i=D_j\), and the rule's condition is
\(x\notin\FV(\delta)\) for every \(\delta\in D_j\).
\begin{itemize}
  \item[(e1)] \(\varphi\in D_j\).  The condition gives
        \(x\notin\FV(\varphi)\).  By the induction hypothesis
        \(\mathcal D_j\) ends in \((\varphi\to\chi_j,E_j)\) with
        \(E_j\subseteq D_j\setminus\{\varphi\}\), so \(x\) is free in no
        member of \(E_j\), and generalization is permitted on that line.
        Let \(\mathcal D_i\) be \(\mathcal D_j\) followed by
        \[
        \begin{array}{ll}
          (\forall x(\varphi\to\chi_j),\ E_j) & \text{Gen},\\
          (\forall x(\varphi\to\chi_j)\to(\varphi\to\forall x\,\chi_j),\
            \varnothing) & \forall2,\\
          (\varphi\to\forall x\,\chi_j,\ E_j) & \text{MP}.
        \end{array}
        \]
        The \(\forall2\) instance satisfies its side condition
        \(x\notin\FV(\varphi)\).  So \(E_i=E_j\subseteq D_i\setminus\{\varphi\}\).
  \item[(e2)] \(\varphi\notin D_j\).  Then
        \(D_j\subseteq(\Gamma\cup\{\varphi\})\setminus\{\varphi\}
        \subseteq\Gamma\).  By \cref{thm:p02-w02-p02-025-support}(2),
        applied to line \(j\) of the given derivation with
        \(\Delta\defeq\Gamma\), some derivation \(\mathcal R\) from
        \(\Gamma\) ends in \((\chi_j,D_j)\).  Let \(\mathcal D_i\) be
        \(\mathcal R\) followed by
        \[
        \begin{array}{ll}
          (\forall x\,\chi_j,\ D_j) & \text{Gen},\\
          (\forall x\,\chi_j\to(\varphi\to\forall x\,\chi_j),\ \varnothing)
            & \text{A1},\\
          (\varphi\to\forall x\,\chi_j,\ D_j) & \text{MP}.
        \end{array}
        \]
        Gen is permitted because its condition concerns only \(D_j\), for
        which it holds in the given derivation.  So
        \(E_i=D_j=D_i\setminus\{\varphi\}\).
\end{itemize}

\emph{Case (f): line \(i\) comes by existential elimination from line
\(j<i\).}  Then \(\chi_j=\theta\to\rho\), \(\chi_i=\exists x\,\theta\to\rho\),
\(D_i=D_j\), and the rule's conditions are \(x\notin\FV(\rho)\) and
\(x\notin\FV(\delta)\) for every \(\delta\in D_j\).
\begin{itemize}
  \item[(f1)] \(\varphi\in D_j\).  The condition gives
        \(x\notin\FV(\varphi)\).  By the induction hypothesis
        \(\mathcal D_j\) ends in \((\varphi\to(\theta\to\rho),E_j)\) with
        \(E_j\subseteq D_j\setminus\{\varphi\}\).  Let \(\mathcal D_i\) be
        \(\mathcal D_j\) followed by:
        \begin{enumerate}
          \item the nineteen lines of \cref{thm:p02-w02-p02-025-permutation}
                for \((\varphi\to(\theta\to\rho))\to(\theta\to(\varphi\to\rho))\)
                and modus ponens, giving
                \((\theta\to(\varphi\to\rho),E_j)\);
          \item existential elimination, giving
                \((\exists x\,\theta\to(\varphi\to\rho),E_j)\);
          \item the nineteen lines of \cref{thm:p02-w02-p02-025-permutation}
                for the instance
                \((\exists x\,\theta\to(\varphi\to\rho))\to
                (\varphi\to(\exists x\,\theta\to\rho))\) and modus ponens,
                giving \((\varphi\to(\exists x\,\theta\to\rho),E_j)\).
        \end{enumerate}
        Step~2 is permitted: its consequent is \(\varphi\to\rho\), and
        \(x\notin\FV(\varphi)\cup\FV(\rho)=\FV(\varphi\to\rho)\); its premise
        line has support \(E_j\subseteq D_j\), in whose members \(x\) is not
        free.  So \(E_i=E_j\subseteq D_i\setminus\{\varphi\}\).
  \item[(f2)] \(\varphi\notin D_j\).  As in (e2),
        \(D_j\subseteq\Gamma\) and \cref{thm:p02-w02-p02-025-support}(2)
        gives a derivation \(\mathcal R\) from \(\Gamma\) ending in
        \((\theta\to\rho,D_j)\).  Let \(\mathcal D_i\) be \(\mathcal R\)
        followed by
        \[
        \begin{array}{ll}
          (\exists x\,\theta\to\rho,\ D_j) & \exists E,\\
          ((\exists x\,\theta\to\rho)\to(\varphi\to(\exists x\,\theta\to\rho)),\
            \varnothing) & \text{A1},\\
          (\varphi\to(\exists x\,\theta\to\rho),\ D_j) & \text{MP}.
        \end{array}
        \]
        Existential elimination is permitted because its conditions concern
        only \(\rho\) and \(D_j\) and hold in the given derivation.  So
        \(E_i=D_j=D_i\setminus\{\varphi\}\).
\end{itemize}

These six cases exhaust the justifications of
\cref{def:p02-hilbert-fol-dependencies}: a premise line has \(\chi_i=\varphi\)
or \(\chi_i\in\Gamma\setminus\{\varphi\}\), and the remaining clauses are
axioms, modus ponens, generalization and existential elimination.  The
induction is complete.  Generalization occurs in \(\mathcal D_i\) only in
cases (e1) and (e2), at a line whose support is \(E_j\) or \(D_j\); in
neither case is it applied to a line whose support contains \(\varphi\).

Finally, the equivalence: clause~1 turns a derivation of \(\psi\) from
\(\Gamma\cup\{\varphi\}\) into one of \(\varphi\to\psi\) from \(\Gamma\), and
clause~2 does the converse.  This is the statement of
\cref{thm:p02-hilbert-fol-deduction}.
\end{proof}
\end{annexitem}

\begin{annexitem}{annex:p02-w02-p02-025-kleene-deduction}{thm:p02-w02-p02-025-kleene-deduction}{Kleene's deduction theorem for the predicate calculus}
\begin{proof}
Throughout, postulate numbers refer to
\cref{def:p02-w02-p02-025-kleene-postulates}, and a deduction comes with an
analysis.  The proof has four steps.

\emph{Step 0: composition.}  If \(\Gamma\derives D_1\), \ldots,
\(\Gamma\derives D_m\) and \(D_1,\ldots,D_m\derives E\), then
\(\Gamma\derives E\).  Take a deduction of \(E\) from \(D_1,\ldots,D_m\) and
replace each member that the analysis justifies as the assumption occurrence
\(D_p\) by a deduction of \(D_p\) from \(\Gamma\).  Each inserted block ends
in \(D_p\) at the position the occurrence held, so every later immediate
consequence still has its premises among the preceding formulas; the members
of the inserted blocks are occurrences from \(\Gamma\), axioms or immediate
consequences of preceding members.  The result is a deduction from
\(\Gamma\), and it ends in \(E\).  The construction applies no rule that the
given deductions do not already apply.

\emph{Step 1: deductions without Rules 9 and 12.}  Suppose
\(A_1,\ldots,A_k\) is a deduction from \(\Gamma,A\) that applies no
rule other than Rule~2.  We show, by course-of-values induction on
\(i\leq k\), that \(\Gamma\derives A\supset A_i\) by a deduction that applies
no rule other than Rule~2.  In each case the displayed members, with
the induction hypothesis and Step~0, form the required deduction.
\begin{itemize}
  \item[(a)] \(A_i\) is an occurrence from \(\Gamma\) or an axiom:
        \(A_i\); \(A_i\supset(A\supset A_i)\) by 1a; \(A\supset A_i\) by
        Rule~2.
  \item[(b)] \(A_i\) is the assumption occurrence \(A\): the deduction
        \[
        \begin{array}{lll}
          1. & A\supset(A\supset A) & \text{1a}\\
          2. & (A\supset(A\supset A))\supset{} & \text{1b}\\
             & \quad((A\supset((A\supset A)\supset A))\supset(A\supset A)) & \\
          3. & (A\supset((A\supset A)\supset A))\supset(A\supset A)
             & \text{Rule 2: }1,2\\
          4. & A\supset((A\supset A)\supset A) & \text{1a}\\
          5. & A\supset A & \text{Rule 2: }4,3
        \end{array}
        \]
        uses no assumption.  Line~1 is 1a with \(B\defeq A\); line~2 is 1b
        with \(B\defeq A\supset A\) and \(C\defeq A\); line~4 is 1a with
        \(B\defeq A\supset A\).
  \item[(d)] \(A_i\) is an immediate consequence by Rule~2 of \(A_j\) and
        \(A_{j'}=A_j\supset A_i\) with \(j,j'<i\).  The induction hypothesis gives
        \(\Gamma\derives A\supset A_j\) and
        \(\Gamma\derives A\supset(A_j\supset A_i)\).  The 1b instance
        \((A\supset A_j)\supset((A\supset(A_j\supset A_i))\supset(A\supset A_i))\)
        and Rule~2 twice give \(A\supset A_i\).
\end{itemize}
These cases exhaust the justifications in a deduction that applies only
Rule~2.  At \(i=k\), \(\Gamma\derives A\supset A_k\).

\emph{Step 2: three derived rules.}  For all formulas \(A,B,C\):
\[
\begin{array}{ll}
 \text{(i)} & A\supset(B\supset C)\derives A\mathbin{\&}B\supset C;\\
 \text{(ii)} & A\mathbin{\&}B\supset C\derives A\supset(B\supset C);\\
 \text{(iii)} & A\supset(B\supset C)\derives B\supset(A\supset C).
\end{array}
\]
For (i), the sequence \(A\mathbin{\&}B\) (assumption);
\(A\mathbin{\&}B\supset A\) (4a); \(A\) (Rule 2); \(A\supset(B\supset C)\)
(assumption); \(B\supset C\) (Rule 2); \(A\mathbin{\&}B\supset B\) (4b);
\(B\) (Rule 2); \(C\) (Rule 2) is a deduction of \(C\) from
\(A\supset(B\supset C),A\mathbin{\&}B\) applying only Rule~2, and Step~1
discharges \(A\mathbin{\&}B\).  For (ii), the sequence \(A\) (assumption);
\(A\supset(B\supset A\mathbin{\&}B)\) (3); \(B\supset A\mathbin{\&}B\)
(Rule 2); \(B\) (assumption); \(A\mathbin{\&}B\) (Rule 2);
\(A\mathbin{\&}B\supset C\) (assumption); \(C\) (Rule 2) is a deduction of
\(C\) from \(A\mathbin{\&}B\supset C,A,B\) applying only Rule~2.  Step~1
discharges \(B\), giving a deduction of \(B\supset C\) from
\(A\mathbin{\&}B\supset C,A\) that again applies only Rule~2, and a second
application discharges \(A\).  For (iii), the sequence \(A\) (assumption);
\(A\supset(B\supset C)\) (assumption); \(B\supset C\) (Rule 2); \(B\)
(assumption); \(C\) (Rule 2) is a deduction of \(C\) from
\(A\supset(B\supset C),B,A\) applying only Rule~2; Step~1 discharges \(A\)
and then \(B\).

\emph{Step 3: the theorem.}  Let \(A_1,\ldots,A_k\) be a deduction of
\(B=A_k\) from \(\Gamma,A\) in which the free variables are held constant for
\(A\).  We show by course-of-values induction on \(i\leq k\) that
\(\Gamma\derives A\supset A_i\).  The initial segment \(A_1,\ldots,A_i\),
with the restricted analysis, is a deduction from \(\Gamma,A\); dependence of
its members on \(A\) is the same as in the whole deduction, because it is
determined by earlier members only; and it applies Rule 9 or 12 only where
the whole deduction does, so the free variables are held constant for \(A\)
in it as well.  The cases are determined by the justification of \(A_i\).

Cases (a), (b) and (d) are the constructions of Step~1, with the present
induction hypothesis.

\emph{Case (e): \(A_i\) is \(C\supset\forall xA(x)\), an immediate
consequence by Rule~9 of \(A_j=C\supset A(x)\), \(j<i\), where \(x\) is not
free in \(C\).}
\begin{itemize}
  \item[(e1)] \(A_j\) depends on \(A\).  Then \(x\) is not free in \(A\):
        otherwise \(x\) would be varied for \(A\), since Rule~9 is applied
        with respect to \(x\) to the premise \(A_j\), which depends on \(A\).
        The induction hypothesis gives
        \(\Gamma\derives A\supset(C\supset A(x))\).  By Step~2(i) and Step~0,
        \(\Gamma\derives A\mathbin{\&}C\supset A(x)\).  A free occurrence of
        \(x\) in \(A\mathbin{\&}C\) would be a free occurrence in \(A\) or in
        \(C\), so \(x\) is not free in \(A\mathbin{\&}C\), and Rule~9 with
        \(A\mathbin{\&}C\) in the place of its \(C\) applies:
        \(\Gamma\derives A\mathbin{\&}C\supset\forall xA(x)\).  By
        Step~2(ii) and Step~0,
        \(\Gamma\derives A\supset(C\supset\forall xA(x))\).
  \item[(e2)] \(A_j\) does not depend on \(A\).  Then neither does \(A_i\),
        whose only premise is \(A_j\).  Let \(\Sigma\) be the subsequence of
        \(A_1,\ldots,A_i\) consisting of the members that do not depend on
        \(A\).  Each of them is an occurrence from \(\Gamma\) (the assumption
        occurrence \(A\) depends on itself), an axiom, or an immediate
        consequence of members that, by the definition of dependence, do
        not depend on \(A\) either and so belong to \(\Sigma\).  Rules~2, 9
        and 12 carry their stipulations on the formulas involved only, not
        on assumption formulas.  So \(\Sigma\) is a deduction from \(\Gamma\)
        ending in \(A_i\).  Append \(A_i\supset(A\supset A_i)\) by 1a and
        \(A\supset A_i\) by Rule~2.
\end{itemize}

\emph{Case (f): \(A_i\) is \(\exists xA(x)\supset C\), an immediate
consequence by Rule~12 of \(A_j=A(x)\supset C\), \(j<i\), where \(x\) is not
free in \(C\).}
\begin{itemize}
  \item[(f1)] \(A_j\) depends on \(A\).  As in (e1), \(x\) is not free in
        \(A\).  The induction hypothesis gives
        \(\Gamma\derives A\supset(A(x)\supset C)\).  By Step~2(iii) and
        Step~0, \(\Gamma\derives A(x)\supset(A\supset C)\).  The variable
        \(x\) is free in neither \(A\) nor \(C\), so it is not free in
        \(A\supset C\), and Rule~12 with \(A\supset C\) in the place of its
        \(C\) gives \(\Gamma\derives\exists xA(x)\supset(A\supset C)\).
        Step~2(iii) and Step~0 give
        \(\Gamma\derives A\supset(\exists xA(x)\supset C)\).
  \item[(f2)] \(A_j\) does not depend on \(A\).  The construction of (e2)
        applies word for word: the members not depending on \(A\) form a
        deduction of \(A_i\) from \(\Gamma\), and 1a with Rule~2 gives
        \(A\supset A_i\).
\end{itemize}

The justifications of a deduction in the predicate calculus are an
assumption occurrence, an axiom, and Rules 2, 9 and 12, so the cases are
exhaustive, and at \(i=k\) we obtain \(\Gamma\derives A\supset B\).

For the converse, a deduction of \(A\supset B\) from \(\Gamma\) followed by
the assumption occurrence \(A\) and \(B\) by Rule~2 is a deduction of \(B\)
from \(\Gamma,A\).
\end{proof}
\end{annexitem}
